\section{The first chamber step}\label{subsec:rg-step-first-vertex}

At $k=0$ write the charge as $q$.
There is no input activity, and $H_0(B)=\bar H_0(B)=1$.
Consequently the exponent in
\eqref{eq:rg-step-fluctuation-exponent} is
\begin{equation}\label{eq:rg-step-first-exponent}
 V_0^x=
  -\langle\bar\Psi^f,(D_{\mathscr T+\mathscr Z}
                                    -D_{\mathscr T})\Psi^f\rangle^{\mathrm{ch}}
  -m\left(\langle\bar\Psi^f,\Psi^f\rangle^{\mathrm{ch}}
                   -\langle\bar\Psi^c,\Psi^c\rangle^{\mathrm{ch}}\right).
\end{equation}
The terms with $\mu$ and $m$ cancel in the connection
difference of the actions; the second difference
of interactions gives the displayed mass term.
Here $\Psi^c=\psi'_L$, $\bar\Psi^c=\bar\psi'_L$,
$\Gamma_0(\mathscr T)=F_0(\mathscr T)^{-1}$,
and $\Psi^f=\Psi^c+\Gamma_0(\mathscr T)W$.
All these are fields in $\mathscr V_0=\Sspace^{\mathrm{ch}}(1,\ell)$.
The identity measure in $W$ has mode measure $1/48$,
and its two-point kernel is $48$ times the identity
matrix of the modes.

For a bond fluctuation $Z$ and an $M$-cube $\square$,
define its local vertex by
\begin{equation}\label{eq:rg-step-first-vertex}
 \begin{split}
 \mathcal V_0(Z,\square)=\frac12\int_\square\dd x\,
   \Big\{&
    \bar\psi(x)^TM_3(x)(e^{iqZ(c,c_3)}-1)
                                  [\psi(x)+\chi(x,x^{(3)})]\\
    &+[\bar\psi(x)+\bar\chi(x,x^{(3)})]^T
               M_3(x^{(3)})(e^{-iqZ(c,c_3)}-1)\psi(x)\Big\},
 \end{split}
\end{equation}
where $x=(c,g)$, $c_3=c+ge_1$ and
$\int_\square\dd x=48^{-1}\sum_{c\in\square}\sum_g$.
The pair $(x,x^{(3)})$ has $d_X=1$ and is admissible
for $\square$ of the second kind. Thus the vertex
has ordinary legs over the cube and precisely the
boundary pair legs allowed in the activity algebra.
Its coefficients are entire functions of $Z$ on
the bonds meeting the cube.

\begin{proposition}\label{prop:rg-step-first-vertex}
Evaluate $\chi=\delta_{\alpha,\mathscr T}f$,
$\bar\chi=\delta_{\alpha,-\mathscr T}\bar f$,
$\psi=f$ and $\bar\psi=\bar f$ at any pair of
antiperiodic odd chamber fields. Then
\begin{equation}\label{eq:rg-step-first-vertex-evaluation}
 \operatorname{ev}\sum_\square\mathcal V_0(Z,\square)
          =-\langle\bar f,(D_{\mathscr T+Z}-D_{\mathscr T})
                                                    f\rangle^{\mathrm{ch}}.
\end{equation}
The two halves have equal images after summing
over cubes. The vertex has the lattice, real-gauge,
charge-conjugation and reality symmetries
\begin{equation}\label{eq:rg-step-first-vertex-symmetries}
 \begin{split}
 a_g\mathcal V_0(gZ,g\square)&=\mathcal V_0(Z,\square),\\
 G_\lambda\mathcal V_0(Z,\square)&=\mathcal V_0(Z,\square),\\
 \mathrm{Cg}\mathcal V_0(-Z,\square)&=\mathcal V_0(Z,\square),\\
 \mathbf K\mathcal V_0(-\overline Z,\square)&=\mathcal V_0(Z,\square).
 \end{split}
\end{equation}
For any generator weights $(w_1,w_2)$,
\begin{equation}\label{eq:rg-step-first-vertex-bound}
 \|\mathcal V_0(Z,\square)\|_{w_1,w_2}
       \le12(1+\rho)M^3w_1(w_1+w_2)
       \sup_{b\text{ oriented and meeting }\square}
                                                      |e^{iqZ(b)}-1|.
\end{equation}
If $q\le1$ and
$q|Z(b)|\le\frac12q^{1-\varepsilon/2}$ on these bonds,
then
\begin{equation}\label{eq:rg-step-first-vertex-small-bound}
 \|\mathcal V_0(Z,\square)\|_{h_0/2}
                 \le6(1+\rho)M^3q^{1/2-\varepsilon/2}
                 \le12M^3q^{1/2-\varepsilon/2}.
\end{equation}
\end{proposition}

\begin{proof}
At a nearest-neighbor pair the six paths reduce to
the bond, so the quotient relations read
\[
 f(x)+\chi(x,x^{(3)})=u_{\mathscr T}(c,c_3)f(x^{(3)}),
 \qquad
 \bar f(x)+\bar\chi(x,x^{(3)})
                         =u_{\mathscr T}(c,c_3)^{-1}\bar f(x^{(3)}).
\]
The first half of the summed vertex is therefore
one half of the integral of
\[
 \bar f(x)^TM_3(x)(e^{iqZ(c,c_3)}-1)
                                  u_{\mathscr T}(c,c_3)f(x^{(3)}).
\]
Reindexing the second half by the measure-preserving
involution $x\mapsto x^{(3)}$ gives the same
integrand: both $Z$ and $\mathscr T$ change sign
on the reversed bond. Only the type-$3$ hop
depends on the connection, and
$u_{\mathscr T+Z}=e^{iqZ}u_{\mathscr T}$.
The difference of the two Dirac operators is
thus the negative of this hopping coefficient,
proving \eqref{eq:rg-step-first-vertex-evaluation}.
Every product contains two antiperiodic fields,
so the reindexing is legitimate on the torus.

For lattice symmetry use
$g^{-1}(x^{(3)})=(g^{-1}x)^{(3)}$ and
$S^{-1}M_3(x)S=M_3(g^{-1}x)$ for the spin
lift $S$. Each half transforms separately, and
reindexing maps $g\square\times O_{\mathrm{cub}}$
onto $\square\times O_{\mathrm{cub}}$ with the
same measure. For a real gauge, the two generators
in each term have the same first photon site and
opposite bar indices, so their phases cancel.
Charge conjugation interchanges the two halves.
Its algebraic rule on the first half gives
\[
 \mathrm{Cg}[\bar\psi(x)^Tm(\psi(x)+\chi(x,x^{(3)}))]
       =(\bar\psi(x)+\bar\chi(x,x^{(3)}))^T
                              (\mathsf C^{-1}m^T\mathsf C)\psi(x).
\]
The sign from exchanging the odd factors cancels
the minus sign in the image of $\bar\psi$.
As
\[
 \mathsf C^{-1}M_3(x)^T\mathsf C=M_3(x^{(3)}),
\]
the first half at $-Z$ becomes the second at $Z$;
the reverse relation gives the other half.
Finally
\[
 \varkappa^T\overline{M_3(x)}\varkappa=M_3(x),\qquad
 \overline{e^{\pm iq(-\overline Z)}-1}=e^{\pm iqZ}-1.
\]
The antilinear map consequently preserves each
half in its reality identity. This proves
\eqref{eq:rg-step-first-vertex-symmetries}.

Put $\delta_x=e^{iqZ(c,c_3)}-1$ and
$\delta'_x=e^{-iqZ(c,c_3)}-1$. Collecting the
ordinary bilinear in \eqref{eq:rg-step-first-vertex},
its coefficient is
$[M_3(x)\delta_x+M_3(x^{(3)})\delta'_x]/96$.
The other two products have coefficients
$M_3(x)\delta_x/96$ and
$M_3(x^{(3)})\delta'_x/96$, and weights
$w_1w_2$.
The ordinary product has weight $w_1^2$.
The chamber sites in the cube and the ordered
pairs with these first sites are distinct, so
these are exactly the coefficients, up to
ordering signs. Each endpoint matrix satisfies
$\sum_{\beta,\beta'}|M_3(\cdot)_{\beta\beta'}|
\le12(1+\rho)$, since
$M_3=3(\rho\mp\gamma_\mu)$ and $\gamma_\mu$
has four entries of modulus one.
There are $48M^3$ chamber sites in the cube.
The triangle inequality for the collected
bilinear therefore gives
\eqref{eq:rg-step-first-vertex-bound}.

Under the smallness condition,
$|e^{iqZ(b)}-1|\le2q|Z(b)|\le q^{1-\varepsilon/2}$.
For $w=h_0/2$,
$w_1(w_1+w_2)=\frac14q^{-1/2}(1+q^\varepsilon)
\le\frac12q^{-1/2}$. Inserting these bounds proves
\eqref{eq:rg-step-first-vertex-small-bound}.
\end{proof}

For $Z=\mathscr Z$ and the evaluation at
$\Psi^f,\bar\Psi^f$, the first term of
\eqref{eq:rg-step-first-exponent} is the image of
the summed vertex. Under the additional charge
condition $96M^3q^{1/4}\le1$, its bound becomes
$\|\mathcal V_0(\mathscr Z,\square)\|_{h_0/2}
\le q^{1/4-5\varepsilon}$ whenever the displayed
bond smallness condition holds. These statements
apply before any expansion of the finite integral.

The first-step density identity retains every
vacuum and determinant factor:
\begin{equation}\label{eq:rg-step-first-density}
 \begin{split}
 (\mathbf T_0^x\mathscr F_0^x[\varepsilon^{\mathrm v},m,0])
                                    (A',\mathscr C';\Psi')\\
  ={}&\det D_1(0)^{-1}
       \exp\Big\{-\langle\bar\Psi',D_1(\mathscr T')\Psi'\rangle_1
                   -Lm\langle\bar\psi',\psi'\rangle\\
      &\hspace{14mm}
                   -L^3(\varepsilon^{\mathrm v}
                        +\delta\varepsilon_0^0)n_1
                   +\sum_X\mathfrak P_0(X,\mathscr C)\Big\}\\
       &\times
              \frac{\det F_0(\mathscr T)}{\det F_0(0)}
        \int\dd\mu_I(W)\,\dd\widehat\mu_0(\zeta)\,e^{V_0^x}.
 \end{split}
\end{equation}
The two determinants of $F_0$ have size $192n_0$.
The original level-zero density, integrated before
the block transformation, is
\begin{equation}\label{eq:rg-step-initial-integrated-density}
 \int\mathrm D\Psi\,\mathscr F_0^x
                     [\varepsilon^{\mathrm v},m,0](A,C;\Psi)
     =e^{-\varepsilon^{\mathrm v}n_0+\sum_X\mathfrak P_0(X,C)}
               \frac{\det(D_{A+C}+\mu+m)}{\det D_0(0)}.
\end{equation}
This follows directly from the weighted determinant
identity. The numerator need not be invertible;
the nonzero denominator uses $\mu>0$.


